\documentclass[../../../include/open-logic-section]{subfiles}

\begin{document}

\olfileid{sth}{spine}{valpha}

\olsection{તબક્કાઓની $V_\alpha$ તરીકે વ્યાખ્યા}

\olref[sth][ordinals][]{chap}માં આપણે સુક્રમો અને (વૉન નોયમનની)
ક્રમવાચક સંખ્યાઓ વ્યાખ્યાયિત કરી. આ પ્રકરણમાં તેનો ઉપયોગ કરીને
ગણોના પદાનુક્રમનું \emph{પોતાનું} સ્વરૂપ સમજશું. તે માટે યાદ કરો કે
\olref[sth][ordinals][opps]{sec}માં આપણે અનુગામી અને સીમા
ક્રમવાચક સંખ્યાઓ વ્યાખ્યાયિત કરી હતી. આ ખ્યાલો આગળની
વ્યાખ્યામાં વાપરીશું:

\begin{defn}\ollabel{defValphas}
	\begin{align*}
	V_\emptyset &\defis \emptyset\\
	V_{\ordsucc{\alpha}} &\defis \Pow{V_\alpha} & & 
	\text{for any ordinal }\alpha\\
	V_{\alpha} &\defis \bigcup_{\gamma < \alpha} V_\gamma & & 
	\text{when }\alpha\text{ is a limit ordinal}
\end{align*}
\end{defn}
\noindent
આ ક્રમવાચક સંખ્યાઓ પર \emph{અનંતાતીત પુનરાવર્તન} વડે કરેલી
વ્યાખ્યા થશે. તેની સરખામણી પ્રાકૃતિક સંખ્યાઓ પરના વિધેયોની
પુનરાવર્તી વ્યાખ્યાઓ સાથે કરવી જોઈએ.\footnote{સરખાવો:
\olref[sfr][infinite][dedekind]{sec}માં સરવાળા, ગુણાકાર અને ઘાતાંકની
વ્યાખ્યાઓ.} પ્રાકૃતિક સંખ્યાઓમાં જેમ, અહીં પણ પાયાનો કિસ્સો અને
અનુગામી કિસ્સાઓ વ્યાખ્યાયિત થાય છે; પરંતુ ક્રમવાચક સંખ્યાઓ માટે
\emph{સીમા} કિસ્સાઓમાં શું થાય છે તે પણ કહેવું પડે છે.

$V_\alpha$ની આ વ્યાખ્યા એક મહત્ત્વનો પડાવ બનશે. આપણે અનૌપચારિક
રીતે ગણોના પદાનુક્રમને \emph{તબક્કાઓ}માં ગણો રચાય છે એમ
સમજાવ્યો હતો. $V_\alpha$ અસરકારક રીતે એ જ તબક્કાઓ છે. જોકે
મહત્ત્વનું એ છે કે આ તબક્કાઓનું \emph{સિદ્ધાંતની અંદરનું}
સ્વરૂપનિરૂપણ છે. સિદ્ધાંતનું સંભવિત \emph{પ્રતિરૂપ} સૂચવવાને
બદલે આપણે આપણા ગણસિદ્ધાંતની \emph{અંદર} જ તબક્કાઓ
વ્યાખ્યાયિત કર્યા હશે.

\end{document}
